\begin{papercard}{habib2019}
\cf{Stages} \stagetag{A} \stagetag{D} \stagetag{W} \stagetag{C} \stagetag{I}. An expert content analysis (heuristic evaluation) of email, advertising and deletion choices on a stratified sample of 150 websites, using a structured Qualtrics template, a double-coded half sample ($\kappa=0.82$), path-length counts and readability scores. \emph{Note:} PDF page 1 is the USENIX cover.

\cf{RQs} (1) Which email, targeted-ad and deletion choices do websites offer? (2) How are those choices presented? \pc{habib2019}{2}{Our empirical content analysis focused on two research questions}. Scope is usability, not legal compliance \pc{habib2019}{4}{Our analysis primarily focuses on usability issues and does not intend to analyze legal compliance}.

\cf{Operational definitions} Email opt-out: any mechanism to stop any email type. Ad opt-out: any mention of an industry opt-out tool or ad-related setting (site, browser, OS). Deletion: any means to delete an account or account-related information, including by email \pc{habib2019}{6}{we defined opt-outs for email communications as mecha}.

\cf{Template and procedure (14 steps)} Visit the homepage and note any cookie notice; create an account with an alias and a study e-mail; check the AdChoices page; read the policy for email, ads, deletion and Do Not Track; then check account settings for each choice. For each choice, record the \emph{location}, the \emph{user actions in the shortest path}, and the information given \pc{habib2019}{6}{For each choice identified, we recorded where the privacy choice is located on the website}. Good and bad practices were noted at each stage of the \emph{Interaction Cycle} \pc{habib2019}{6}{At every stage, researchers also made note of practices for}.

\cf{Pilot} Six rounds of pilot testing on 25 sites from the Alexa top 50. Two researchers coded independently, reconciled, and revised the template questions for shared understanding. Pilot data were excluded \pc{habib2019}{6}{To refine the template, our research team conducted six rounds of pilot testing}.

\cf{Sample} 150 sites from the Alexa global top 10,000 (22 March 2018), stratified by reach into top (1--200), middle (201--5,000) and bottom ($>$5,000) tiers, 50 each, with thresholds where reach levelled off \pc{habib2019}{6}{These thresholds were identified by plotting websites’ reach against their rank}. Exclusions: porn sites, sites needing non-email verification or sensitive data such as SSN or credit card to register, and non-English sites \pc{habib2019}{6}{In constructing our sample, we ex}. Analysed April--October 2018. The 37 sites that updated policies around the GDPR date were re-analysed, and pre/post versions compared \pc{habib2019}{6}{In October 2018 we reviewed all websites in our dataset that had been analyzed prior to May 25}.

\cf{Reliability} Coders were trained on several sites. Two researchers independently coded 75 sites (50\,\%) spread across tiers. Cohen's $\kappa=0.82$ was \emph{averaged over the presence/absence questions}, disagreements were reconciled, and the other 75 were single-coded \pc{habib2019}{7}{Cohen’s Kappa (κ = 0.82) was averaged over the questions}. Each site took 5--58 minutes (mean 21). Data were collected in a private Chrome window, fresh for each site \pc{habib2019}{7}{researchers collected data in Google Chrome’s private browsing mode}.

\cf{Measures} \emph{Readability} of choice-related policy text (readable.io): FRES about 40 and FKGL about 13, i.e.\ university level, with over 90\,\% ``fairly difficult'' or worse \pc{habib2019}{7}{The analyzed text for all three types of privacy choices on the Flesch-Kincaid Grade Level}. \emph{Section headings}: n-grams counted; no n-gram appeared in more than 20 policies \pc{habib2019}{10}{However, no single n-gram occurred in more than 20 of the policies we}. \emph{Path length}: clicks, hovers, form fields, radio buttons and checkboxes from the homepage up to (not including) the apply action \pc{habib2019}{10}{We counted user actions as the number of clicks, hovers, form fields, radio buttons}. The New York Times deletion form required 38 actions \pc{habib2019}{10}{38 user actions were required to complete the New York Times’ data deletion request form}.

\cf{Findings} Choices were prevalent: 89\,\% of applicable sites offered email and ad opt-outs and 74\,\% offered deletion \pc{habib2019}{8}{89\% of websites with email marketing or targeted advertising offered opt-outs}. But only 27 policies linked directly to a deletion tool, and 81 gave e-mail instructions \pc{habib2019}{9}{only 27 policies included a direct link to a data deletion tool or request form}. 27\,\% of DAA/NAI links went to homepages instead of the opt-out tool \pc{habib2019}{9}{27\% of opt-out links pointing to the DAA or NAI directed visi}. 90 of 108 sites with deletion gave no deletion time frame \pc{habib2019}{10}{Ninety of 108 websites offering deletion did not describe a time frame}. Missing or misleading policy information, broken links and non-functioning mechanisms appeared on 71 sites \pc{habib2019}{11}{We also recorded 15 instances in which provided links to rel}. Only 5\,\% honoured Do Not Track.

\cf{Limitations stated} English-only, Alexa top 10k, exclusions of sites needing sensitive data, US IP addresses (EU-specific choices unobserved), and no evidence about actual user comprehension \pc{habib2019}{7}{Additionally, since our analysis was conducted using US IP addresses}.

\cf{Recommendations by interaction stage} \emph{Finding}: standardised headings and a central location. \emph{Learning}: a single mechanism per choice, and relevant umbrella policies. \emph{Using}: one-click opt-outs, auto-save to avoid post-completion errors, actionable links. \emph{Understanding}: describe what a choice does, its scope and its time frame \pc{habib2019}{13}{Removing this extra step would avoid post-completion errors}.

\cf{Use for this thesis} Probably the best methodological template for an expert walkthrough of privacy choices. (1) Build a \emph{fixed step-by-step capture template} (home, store page, policy, settings) and pilot it in rounds with two coders before data collection. (2) Use \emph{study-only accounts} and a clean browser or device for each app. (3) Record location, shortest-path action counts (counting rules fixed in advance, apply action excluded) and readability. (4) Double-code a sizeable subset and report $\kappa$ \emph{per question type}; Habib averaged only over presence items, a limitation to avoid or state. (5) Its US-IP caveat is the mirror image of a German-only sample, so record the region and store. (6) Its four-stage Interaction Cycle (finding, learning, using, understanding) can structure walkthrough results.
\end{papercard}
